\section{CAF--Immune Cell Crosstalk in Breast Cancer}
\label{sec:immune}

The bidirectional communication between CAFs and immune cells is a central feature of the breast tumour microenvironment. CAFs shape the immune landscape through multiple mechanisms, while immune cells reciprocally regulate CAF activation and function. This crosstalk has profound implications for immunotherapy efficacy.

\subsection{CAFs as immune regulators}

CAFs modulate both innate and adaptive immune responses through paracrine signalling, direct cell--cell contact, and ECM remodelling.

\subsubsection{Suppression of anti-tumour immunity}

Multiple CAF subpopulations contribute to immune suppression through distinct mechanisms:

\textbf{iCAF-mediated immunosuppression:} iCAFs secrete a cocktail of immunosuppressive factors including IL-6, CXCL12, CCL2, and prostaglandin E2 (PGE2) \citep{erve2024}. IL-6 activates STAT3 signalling in CD8+ T cells, impairing their cytotoxic function and promoting exhaustion. CXCL12 recruits regulatory T cells (Tregs) and myeloid-derived suppressor cells (MDSCs), creating an immunosuppressive milieu. CCL2 attracts monocytes that differentiate into tumour-associated macrophages (TAMs) with M2 polarization, further suppressing anti-tumour immunity.

\textbf{myCAF physical barrier:} myCAFs produce a dense, collagen-rich ECM that physically excludes immune cells from tumour cell clusters \citep{chhabra2023}. This immune-excluded phenotype---characterized by immune cells trapped in the stroma rather than infiltrating tumour nests---is associated with resistance to immune checkpoint inhibitors (ICIs). The mechanical barrier created by myCAFs is reinforced by their contractility, which compresses lymphatic vessels and reduces immune cell trafficking \citep{fan2023}.

\textbf{apCAF tolerogenic antigen presentation:} apCAFs present tumour antigens to CD4+ T cells via MHC class II without providing adequate co-stimulation, inducing T cell anergy and tolerance \citep{erve2024}. This mechanism is analogous to peripheral tolerance induction by tolerogenic dendritic cells and may contribute to the failure of anti-tumour immune responses.

\subsubsection{Promotion of pro-tumorigenic immunity}

Paradoxically, some CAF populations promote inflammatory responses that support tumour growth:

\textbf{Neutrophil recruitment:} iCAF-secreted CXCL1 and CXCL8 attract neutrophils to the tumour site \citep{acharyya2012}. These tumour-associated neutrophils (TANs) can adopt a pro-tumorigenic N2 phenotype, producing matrix metalloproteinases (MMPs) that degrade the basement membrane and facilitate invasion. TANs also secrete reactive oxygen species (ROS) that promote DNA damage and genomic instability in tumour cells.

\textbf{Th17 polarization:} iCAF-derived IL-6, combined with TGF-$\beta$, drives CD4+ T cell differentiation toward the Th17 lineage \citep{chang2014}. Th17 cells produce IL-17, which promotes tumour angiogenesis and recruits immunosuppressive myeloid cells. In breast cancer, high Th17 infiltration is associated with aggressive disease and poor prognosis.

\textbf{Macrophage polarization:} CAFs promote the differentiation of monocytes into M2-polarized TAMs through secretion of CSF-1, IL-10, and TGF-$\beta$ \citep{qian2011}. M2 TAMs suppress anti-tumour immunity, promote angiogenesis, and facilitate metastasis, creating a feed-forward loop that reinforces the immunosuppressive microenvironment.

\subsection{Immune regulation of CAFs}

The crosstalk is bidirectional: immune cells also regulate CAF activation and function.

\textbf{TGF-$\beta$ from regulatory T cells:} Tregs secrete high levels of TGF-$\beta$, which activates fibroblasts and promotes their differentiation into myCAFs \citep{donnelly2015}. This creates an immunosuppressive niche where Tregs and myCAFs mutually reinforce each other's functions.

\textbf{Macrophage-derived IL-1$\beta$:} M1-polarized macrophages secrete IL-1$\beta$, which activates NF-$\kappa$B signalling in fibroblasts, promoting their differentiation into iCAFs \citep{erve2024}. This mechanism links innate immune activation to CAF heterogeneity, with the inflammatory state of the myeloid compartment determining the balance between myCAF and iCAF subpopulations.

\textbf{T cell-derived IFN-$\gamma$:} Activated CD8+ T cells secrete IFN-$\gamma$, which can reprogram CAFs toward an anti-tumorigenic phenotype \citep{kato2018}. IFN-$\gamma$-treated CAFs upregulate MHC class I and present tumour antigens, potentially enhancing anti-tumour immunity. However, chronic IFN-$\gamma$ exposure can also induce CAF senescence and immunosuppressive factor production, highlighting the context-dependent nature of immune--CAF interactions.

\subsection{Implications for immunotherapy}

The CAF--immune crosstalk has direct implications for the efficacy of immunotherapies in breast cancer:

\textbf{Resistance to immune checkpoint inhibitors:} The immune-excluded phenotype created by myCAF capsules is a major mechanism of resistance to anti-PD-1/PD-L1 therapy \citep{chhabra2023}. Immune cells cannot reach tumour cells to exert their cytotoxic effects, rendering checkpoint blockade ineffective. This explains why breast cancers with high stromal content---particularly TNBC with dense myCAF infiltration---show poor response to immunotherapy despite high tumour mutational burden.

\textbf{Enhancement of CAR-T cell therapy:} CAF-targeted strategies can enhance the efficacy of chimeric antigen receptor T (CAR-T) cell therapy. CAR-T cells engineered to secrete FAP-targeted antibodies can eliminate FAP+ CAFs, disrupting the immunosuppressive niche and improving T cell infiltration \citep{wang2014}.

\textbf{Combination strategies:} The CAF--immune crosstalk suggests that optimal immunotherapy in breast cancer may require combination approaches that simultaneously target CAFs and immune checkpoints. For example, anti-TGF-$\beta$ therapy can reduce myCAF activation while enhancing T cell infiltration, potentially synergizing with anti-PD-1 therapy \citep{chhabra2023}.

\subsection{CAF-derived extracellular vesicles in immune modulation}

CAFs release extracellular vesicles (EVs)---including exosomes and microvesicles---that carry bioactive cargo (proteins, lipids, miRNAs) to immune cells \citep{zhao2023}. These EVs represent a novel mechanism of CAF--immune communication:

\textbf{miRNA transfer:} CAF-derived exosomes carry miRNAs (miR-21, miR-155, miR-200 family) that reprogram recipient immune cells. miR-21 promotes M2 macrophage polarization, while miR-155 enhances T cell exhaustion \citep{zhao2023}.

\textbf{Metabolic reprogramming:} CAF-derived EVs carry metabolic enzymes and metabolites (lactate, glutamine) that alter the metabolic state of immune cells, shifting them from oxidative phosphorylation to aerobic glycolysis. This metabolic reprogramming impairs T cell effector function and promotes immune suppression \citep{zhao2023}.

\textbf{Antigen transfer:} CAF-derived EVs can carry tumour antigens, potentially enabling cross-presentation by dendritic cells and enhancing anti-tumour immunity. However, the net effect of CAF-derived EVs on immune function appears to be immunosuppressive in most contexts, though this remains an active area of investigation.

\subsection{Therapeutic targeting of CAF--immune crosstalk}

Understanding the CAF--immune crosstalk has revealed multiple therapeutic opportunities:

\textbf{Blocking CAF-derived immunosuppressive factors:} Antibodies or small molecules targeting IL-6, CXCL12, or CCL2 can disrupt the immunosuppressive signalling network and enhance anti-tumour immunity \citep{serve2023}. Clinical trials of IL-6 receptor antibodies (tocilizumab) in combination with ICIs are underway in breast cancer.

\textbf{Reprogramming CAFs to enhance immunogenicity:} Agents that shift CAFs from a suppressive to an immunostimulatory phenotype---such as all-trans retinoic acid (ATRA) or vitamin D---can enhance the efficacy of immunotherapy \citep{erve2024}. These approaches aim to convert immune-excluded tumours into immune-inflamed tumours.

\textbf{Disrupting the physical barrier:} Strategies that degrade the myCAF-produced ECM---such as hyaluronidase (PEGPH20) or collagenase---can enhance immune cell infiltration into tumour nests \citep{provenzano2012}. However, these approaches must be carefully balanced to avoid promoting tumour invasion through basement membrane degradation.
